\chapter{A GENERAL VARIATIONAL PRINCIPLE}
\label{ch:14}

\section{Introduction}\label{sec:121}

In the preceding chapters we have considered instabilities of hydrodynamic and hydromagnetic systems arising from many causes: from adverse distributions of temperature and density, of angular velocity and angular momentum, and from shear, and from gravity and capillarity. In this chapter we shall enunciate a general variational principle which can well form the starting point for further developments.

\section{A variational principle for treating the stability of hydromagnetic systems}\label{sec:122}

To be specific, suppose that the initial state is one in which an incompressible fluid of infinite electrical conductivity is confined in a volume $V$ bounded by a surface $S$. Let the space surrounding the fluid, either finite or infinite, be a vacuum. Suppose further that in the stationary state a magnetic field $\mathbf{H}$ prevails.

We shall distinguish the variables pertaining to the interior and to the exterior of the volume $V$ by superscripts `in' and `ex', if there should be an ambiguity in the context as to which is meant. When there is no such ambiguity we shall dispense with the distinguishing superscripts.

In the stationary state the equations which must be satisfied inside $V$ are
\begin{equation}
\label{eq:14-1}
\frac{\partial \Pi}{\partial x_i} = \frac{1}{4\pi} \frac{\partial}{\partial x_k} H_k H_i,
\end{equation}
where
\begin{equation}
\label{eq:14-2}
\Pi = p + |\mathbf{H}|^2 / 8\pi.
\end{equation}
Outside of $V$, the field, being that of a vacuum, must be current free; accordingly,
\begin{equation}
\label{eq:14-3}
\operatorname{curl} \mathbf{H}^{(\mathrm{ex})} = 0.
\end{equation}
On the bounding surface $S$ the continuity of the normal stress requires
\begin{equation}
\label{eq:14-4}
p^{(\mathrm{in})} + \frac{|\mathbf{H}^{(\mathrm{in})}|^2}{8\pi} = \frac{|\mathbf{H}^{(\mathrm{ex})}|^2}{8\pi}.
\end{equation}
Also, we shall restrict our present considerations to the case when
\begin{equation}
\label{eq:14-5}
\mathbf{H} \cdot d\mathbf{S} = 0 \quad \text{everywhere on } S.
\end{equation}
This condition states that the lines of force of the magnetic field lie on the surface: a necessary condition if infinite accelerations on the surface are \emph{not} to arise.

The equations governing a first order perturbation of the initial state are
\begin{equation}
\label{eq:14-6}
\rho \frac{\partial^2 u_i}{\partial t^2} = -\frac{\partial \varpi}{\partial x_i} + \frac{1}{4\pi} \frac{\partial}{\partial x_k} (H_j h_k + h_j H_k)
\end{equation}
and
\begin{equation}
\label{eq:14-7}
\frac{\partial h_i}{\partial t} = \frac{\partial}{\partial x_k} (H_k u_i - u_k H_i),
\end{equation}
where $\mathbf{h}$ denotes the perturbation in the field and
\begin{equation}
\label{eq:14-8}
\varpi = \delta p + \frac{1}{4\pi} H_j h_j.
\end{equation}
It is convenient to define the displacement $\boldsymbol{\xi}$ by the equation
\begin{equation}
\label{eq:14-9}
u_i = \frac{\partial \xi_i}{\partial t},
\end{equation}
so that equation~(\ref{eq:14-7}) immediately integrates to give
\begin{equation}
\label{eq:14-10}
h_i = \frac{\partial}{\partial x_k} (H_k \xi_i - \xi_k H_i).
\end{equation}
For an incompressible fluid $\boldsymbol{\xi}$ must, of course, be solenoidal.

Suppose now that all the quantities describing the perturbed state have the time dependence $e^{i\sigma t}$. Then, the equation of motion~(\ref{eq:14-6}) gives
\begin{equation}
\label{eq:14-11}
\sigma^2 \rho \xi_i = \frac{\partial \varpi}{\partial x_i} - \frac{1}{4\pi} \frac{\partial}{\partial x_k} (H_j h_k + h_j H_k).
\end{equation}
Also, the perturbation in the field outside $V$ (i.e.\ in $V^{(\mathrm{ex})}$) must satisfy the equations
\begin{equation}
\label{eq:14-12}
\operatorname{curl} \mathbf{h}^{(\mathrm{ex})} = 0 \quad \text{and} \quad \operatorname{curl} \mathbf{E}^{(\mathrm{ex})} = -i\sigma \mathbf{h}^{(\mathrm{ex})}.
\end{equation}
Solutions of equations~(\ref{eq:14-10})--(\ref{eq:14-12}) must be sought which satisfy certain boundary conditions. These conditions which, strictly, must be satisfied on the displaced boundary, $\mathbf{S} + \delta \mathbf{S}$, are
\begin{equation}
\label{eq:14-13}
\Delta(p + |\mathbf{H}|^2/8\pi) = 0,\footnotemark
\end{equation}
\footnotetext{In equations~(\ref{eq:14-13})--(\ref{eq:14-16}), $\mathbf{H}$ stands for the total field, the perturbed plus the unperturbed.}
\begin{equation}
\label{eq:14-14}
\mathbf{N} \cdot \mathbf{H} = 0,
\end{equation}
and
\begin{equation}
\label{eq:14-15}
\mathbf{N} \times \Delta(\mathbf{E}) = (\mathbf{N} \cdot \mathbf{u}) \Delta(\mathbf{H}) = i\sigma (\mathbf{N} \cdot \boldsymbol{\xi}) \Delta(\mathbf{H}),
\end{equation}
where $\Delta(f)$ is the jump a quantity $f$ experiences on $S + \delta S$, and $\mathbf{N}$ is the unit outward normal to this surface. The first of these conditions ensures the continuity of the normal stress, the second, the vanishing of the normal component of $\mathbf{H}$ (in accordance with the assumption~(\ref{eq:14-5})), and the third derives from the relation
\begin{equation}
\label{eq:14-16}
\mathbf{E} + \mathbf{u} \times \mathbf{H} = 0.
\end{equation}
In applying the conditions~(\ref{eq:14-13})--(\ref{eq:14-15}) one reduces them to conditions on $S$ by expanding all the quantities, consistently, to the first order.

Equations~(\ref{eq:14-11}) and~(\ref{eq:14-12}) together with the boundary conditions~(\ref{eq:14-13})--(\ref{eq:14-15}) constitute a characteristic value problem for $\sigma^2$. If $\sigma^{(n)}$ is a characteristic value, we shall distinguish the proper solutions belonging to it by the same superscript.

Consider equation~(\ref{eq:14-11}) belonging to $\sigma^{(n)}$ and after multiplication by $\xi_i^{(m)}$ (belonging to a different characteristic value $\sigma^{(m)}$) integrate over the volume $V$ occupied by the fluid. We obtain
\begin{equation}
\label{eq:14-17}
(\sigma^{(n)})^2 \int_{V^{(\mathrm{in})}} \rho \xi_i^{(n)} \xi_i^{(m)} \, dV - \int_V \frac{\partial}{\partial x_i} (\varpi^{(n)} \xi_i^{(m)}) \, dV - \frac{1}{4\pi} \int_V H_j \frac{\partial h_k^{(n)}}{\partial x_k} \xi_i^{(m)} \, dV - \frac{1}{4\pi} \int_V h_j^{(n)} \frac{\partial H_k}{\partial x_k} \xi_i^{(m)} \, dV.
\end{equation}

Integrating by parts the first two integrals on the right-hand side of~(\ref{eq:14-17}), we obtain
\begin{equation}
\label{eq:14-18}
(\sigma^{(n)})^2 \int \rho \xi_i^{(n)} \xi_i^{(m)} \, dV = \int \varpi^{(n)} \xi_i^{(m)} \, dS_i + \frac{1}{4\pi} \int \left( H_k h_i^{(n)} \frac{\partial \xi_i^{(m)}}{\partial x_k} - h_j^{(n)} \xi_i^{(m)} \frac{\partial H_j}{\partial x_i} \right) dV,
\end{equation}
the integrated part from the second integral vanishing on account of~(\ref{eq:14-5}). Now inserting for $\mathbf{h}$ from equation~(\ref{eq:14-10}), we can transform the volume integral on the right-hand side of equation~(\ref{eq:14-18}) as follows:
\begin{align}
\label{eq:14-19}
& \int H_k \left(H_j \frac{\partial \xi_i^{(n)}}{\partial x_j} - \xi_j^{(n)} \frac{\partial H_i}{\partial x_j}\right) \frac{\partial \xi_i^{(m)}}{\partial x_k} \, dV - \int \frac{\partial}{\partial x_j}\left(H_j \xi_i^{(n)} - \xi_j^{(n)} H_i\right) H_k \frac{\partial \xi_i^{(m)}}{\partial x_k} \, dV \notag \\[4pt]
& \quad = \int \left(H_k \frac{\partial \xi_i^{(n)}}{\partial x_j}\right) \left(H_j \frac{\partial \xi_i^{(m)}}{\partial x_k}\right) dV + \int \frac{\partial}{\partial x_k} \left(\xi_j^{(n)} \frac{\partial H_j}{\partial x_j} + H_j \frac{\partial \xi_j^{(n)}}{\partial x_j}\right) \xi_i^{(m)} \, dV \notag \\[4pt]
& \qquad - \frac{1}{4\pi} \left( H_j \frac{\partial H_i}{\partial x_j} \xi_i^{(m)} + H_i \frac{\partial H_j}{\partial x_j} \xi_i^{(m)} \right) dV.
\end{align}

Interchanging the suffixes $j$ and $k$ in the second term inside the integral on the last line of~(\ref{eq:14-19}), we obtain
\begin{equation}
\label{eq:14-20}
-\int H_j \frac{\partial H_k}{\partial x_j} \frac{\partial}{\partial x_k} (\xi_i^{(n)} \xi_i^{(m)}) \, dV = \int \xi_i^{(n)} \xi_i^{(m)} H_j \frac{\partial^2 H_k}{\partial x_j \partial x_k} \, dV.
\end{equation}

Combining this with the second integral on the second line of~(\ref{eq:14-19}), we have
\begin{equation}
\label{eq:14-21}
\int \xi_i^{(n)} \xi_i^{(m)} \left( \frac{\partial H_j}{\partial x_k} \frac{\partial H_j}{\partial x_k} + H_j \frac{\partial^2 H_k}{\partial x_j \partial x_k} \right) dV = 4\pi \int \xi_i^{(n)} \xi_i^{(m)} \frac{\partial^2 \Pi}{\partial x_k \partial x_k} \, dV,
\end{equation}
where, in the last step, we have made use of equation~(\ref{eq:14-1}). Thus equation~(\ref{eq:14-18}) now becomes
\begin{align}
\label{eq:14-22}
(\sigma^{(n)})^2 \int \rho \xi_i^{(n)} \xi_i^{(m)} \, dV &= \int \varpi^{(n)} \xi_i^{(m)} \, dS_i \notag \\
&\quad + \frac{1}{4\pi} \int \left( H_k \frac{\partial \xi_i^{(n)}}{\partial x_j} H_j \frac{\partial \xi_i^{(m)}}{\partial x_k} + 4\pi \xi_i^{(n)} \xi_i^{(m)} \frac{\partial^2 \Pi}{\partial x_k \partial x_k} \right) dV.
\end{align}

It remains to consider the surface integral in this last equation. We now make use of the condition which follows from~(\ref{eq:14-13}) requiring the continuity of the total pressure on $\mathbf{S} + \delta \mathbf{S}$; we have (cf.\ equation XII (369))
\begin{equation}
\label{eq:14-23}
(\varpi^{(n)})_{\mathrm{in}} = -\frac{1}{4\pi} \mathbf{H}^{(\mathrm{ex})} \cdot \mathbf{h}^{(\mathrm{ex};n)} + (\mathbf{N}^{(0)} \cdot \boldsymbol{\xi}^{(n)})[\mathbf{N}^{(0)} \cdot \operatorname{grad} \Delta_S(\Pi)],
\end{equation}
where $\Delta_S(f)$ now denotes the jump which a quantity $f$ experiences on $S$, and $\mathbf{N}^{(0)}$ is the unit outward normal to $S$. Substituting this expression for $(\varpi^{(n)})_{\mathrm{in}}$ in the surface integral occurring in equation~(\ref{eq:14-22}), we have
\begin{align}
\label{eq:14-24}
\int \varpi^{(n)} \xi_i^{(m)} \, dS_i &= \int \left(\mathbf{N}^{(0)} \cdot \boldsymbol{\xi}^{(n)}\right) (\mathbf{N}^{(0)} \cdot \boldsymbol{\xi}^{(m)}) [\mathbf{N}^{(0)} \cdot \operatorname{grad}\, \Delta_S(\Pi)] \, dS \notag \\
&\quad + \frac{1}{4\pi} \int \left(\mathbf{H}^{(\mathrm{ex})} \cdot \mathbf{h}^{(\mathrm{ex};n)}\right) (\mathbf{N}^{(0)} \cdot \boldsymbol{\xi}^{(m)}) \, dS.
\end{align}
Now from equation~(\ref{eq:14-15}) it follows that
\begin{equation}
\label{eq:14-25}
\mathbf{N}^{(0)} \times \mathbf{E}^{(\mathrm{ex};n)} = i\sigma^{(n)} (\mathbf{N}^{(0)} \cdot \boldsymbol{\xi}^{(n)}) \mathbf{H}^{(\mathrm{ex})}.
\end{equation}
The second integral on the right-hand side of equation~(\ref{eq:14-24}) thus becomes
\begin{equation}
\label{eq:14-26}
\frac{1}{4\pi i \sigma^{(n)}} \int \mathbf{h}^{(\mathrm{ex};n)} \cdot (\mathbf{N}^{(0)} \times \mathbf{E}^{(\mathrm{ex};m)}) \, dS = -\frac{1}{4\pi i \sigma^{(n)}} \int \epsilon_{ijk} \, dS_i h_j^{(\mathrm{ex};n)} E_k^{(\mathrm{ex};m)}.
\end{equation}
The surface integral can now be transformed by Gauss's theorem into a volume integral over $V^{(\mathrm{ex})}$ exterior to $V$; thus,
\begin{equation}
\label{eq:14-27}
\frac{1}{4\pi i \sigma^{(n)}} \int_{V^{(\mathrm{ex})}} \epsilon_{ijk} \frac{\partial}{\partial x_i} (h_j^{(n)} E_k^{(m)}) \, dV
= -\frac{1}{4\pi} \int_{V^{(\mathrm{ex})}} \mathbf{h}^{(n)} \cdot \operatorname{curl} \mathbf{E}^{(m)} \, dV = \frac{1}{4\pi} \int_{V^{(\mathrm{ex})}} h_i^{(n)} h_i^{(m)} \, dV.
\end{equation}

The final form of equation~(\ref{eq:14-22}) is, therefore,
\begin{align}
\label{eq:14-28}
(\sigma^{(n)})^2 \int_{V^{(\mathrm{in})}} \rho \xi_i^{(n)} \xi_i^{(m)} \, dV &= \int_S (\mathbf{N}^{(0)} \cdot \boldsymbol{\xi}^{(n)})(\mathbf{N}^{(0)} \cdot \boldsymbol{\xi}^{(m)})[\mathbf{N}^{(0)} \cdot \operatorname{grad}\, \Delta_S(\Pi)] \, dS \notag \\
&\quad + \frac{1}{4\pi} \int \left( H_k \frac{\partial \xi_i^{(n)}}{\partial x_j} H_j \frac{\partial \xi_i^{(m)}}{\partial x_k} + 4\pi \xi_i^{(n)} \xi_i^{(m)} \frac{\partial^2 \Pi}{\partial x_k \partial x_k} \right) dV + \frac{1}{4\pi} \int_{V^{(\mathrm{ex})}} h_i^{(n)} h_i^{(m)} \, dV.
\end{align}

We observe that the expression on the right-hand side of equation~(\ref{eq:14-28}) is symmetrical in $\lambda$ and $\mu$. Therefore
\begin{equation}
\label{eq:14-29}
\int_{V^{(\mathrm{in})}} \rho \xi_i^{(n)} \xi_i^{(m)} \, dV = 0 \quad (\lambda \neq \mu).
\end{equation}

This clearly establishes the self-adjoint character of the underlying characteristic value problem. Also, by writing in equation~(\ref{eq:14-28}) the complex conjugates of $\xi_i^{(n)}$ and $h_i^{(n)}$ in place of $\xi_i^{(m)}$ and $h_i^{(m)}$ and suppressing the superscripts, we obtain
\begin{align}
\label{eq:14-30}
\sigma^2 \int_{V^{(\mathrm{in})}} \rho |\boldsymbol{\xi}|^2 \, dV &= \int_S |\mathbf{N}^{(0)} \cdot \boldsymbol{\xi}|^2 [\mathbf{N}^{(0)} \cdot \operatorname{grad}\, \Delta_S(\Pi)] \, dS \notag \\
&\quad + \frac{1}{4\pi} \int \left[ \left| H_j \frac{\partial \xi_i}{\partial x_j} \right|^2 + 4\pi \xi_i \xi_i^* \frac{\partial^2 \Pi}{\partial x_k \partial x_k} \right] dV + \frac{1}{4\pi} \int_{V^{(\mathrm{ex})}} |\mathbf{h}|^2 \, dV = \Sigma \;\text{(say)}.
\end{align}

All the quantities on the right-hand side of this equation are clearly real. Therefore $\sigma^2$ must be real and can be either positive or negative. In the former case we will have stability; and in the latter case, instability:
\begin{equation}
\label{eq:14-31}
\Sigma > 0 \quad \text{implies stability} \qquad \text{and} \qquad \Sigma < 0 \quad \text{implies instability.}
\end{equation}

From the self-adjoint character of this problem, it is apparent that equation~(\ref{eq:14-30}) \emph{provides the basis for a variational treatment of the problem}.

\subsection*{(a) The effect of viscosity}

The preceding analysis can be readily modified to allow for effects of viscosity. The modification that is necessary consists only in replacing equation~(\ref{eq:14-11}) by (see Chapter~II, \S\,7\,(b), equations (15)--(17))
\begin{equation}
\label{eq:14-32}
\sigma^2 \rho \xi_i = \frac{\partial \varpi_{ij}}{\partial x_j} - \frac{1}{4\pi} \frac{\partial}{\partial x_k} (H_j h_k + h_j H_k),
\end{equation}
where
\begin{equation}
\label{eq:14-33}
\varpi_{ij} = \left(\delta p + \frac{H_j h_j}{4\pi}\right) \delta_{ij} - P_{ij}
\end{equation}
and
\begin{equation}
\label{eq:14-34}
P_{ij} = \rho \nu \left( \frac{\partial u_i}{\partial x_j} + \frac{\partial u_j}{\partial x_i} \right)
\end{equation}
is the viscous stress tensor. The equation governing the magnetic field is unchanged; but the boundary condition~(\ref{eq:14-13}) must now be replaced by
\begin{equation}
\label{eq:14-35}
N_i \Delta(P_{ij}) = 0.
\end{equation}

It is easy to trace the effect of replacing the scalar pressure $\varpi$ by the tensor pressure $\varpi_{ij}$ in the preceding analysis: in place of the simple surface integral on the right-hand side of equation~(\ref{eq:14-18}) we now have
\begin{equation}
\label{eq:14-36}
\int \varpi_{ij}^{(n)} \xi_i^{(m)} \, dS_j + \int P_{ij}^{(n)} \frac{\partial \xi_i^{(m)}}{\partial x_j} \, dV.
\end{equation}
By virtue of the changed boundary condition~(\ref{eq:14-35}), the further reduction of the surface integral in~(\ref{eq:14-36}) proceeds exactly as before and leads to the same final result. The net change in equation~(\ref{eq:14-28}) is then the addition of the volume integral
\begin{equation}
\label{eq:14-37}
\int P_{ij}^{(n)} \frac{\partial \xi_i^{(m)}}{\partial x_j} \, dV = \rho \nu (i\sigma^{(n)}) \int \left( \frac{\partial \xi_i^{(n)}}{\partial x_j} + \frac{\partial \xi_j^{(n)}}{\partial x_i} \right) \frac{\partial \xi_i^{(m)}}{\partial x_j} \, dV
\end{equation}
to the terms on the right-hand side. The corresponding form of equation~(\ref{eq:14-30}) is therefore
\begin{equation}
\label{eq:14-38}
\sigma^2 \int_{V^{(\mathrm{in})}} \rho |\boldsymbol{\xi}|^2 \, dV = \Sigma + \rho \nu (i\sigma) \int \left( \frac{\partial \xi_i}{\partial x_j} + \frac{\partial \xi_j}{\partial x_i} \right) \frac{\partial \xi_i^*}{\partial x_j} \, dV.
\end{equation}

The volume integral on the right-hand side represents the dissipation due to viscosity (cf.\ equation~II\,(29)). Thus letting
\begin{equation}
\label{eq:14-39}
\Phi = \tfrac{1}{2} \rho \nu \int \left| \frac{\partial \xi_i}{\partial x_j} + \frac{\partial \xi_j}{\partial x_i} \right|^2 dV
\end{equation}
and
\begin{equation}
\label{eq:14-40}
I = \int \rho |\boldsymbol{\xi}|^2 \, dV
\end{equation}
we can rewrite equation~(\ref{eq:14-38}) in the form
\begin{equation}
\label{eq:14-41}
(i\sigma)^2 I + (i\sigma) \Phi + \Sigma = 0.
\end{equation}
The roots of this equation are
\begin{equation}
\label{eq:14-42}
i\sigma = \frac{1}{2I} \left\{ -\Phi \pm \sqrt{\Phi^2 - 4\Sigma I} \right\}.
\end{equation}

From this equation it follows that \emph{a necessary and sufficient condition for stability is that $\Sigma > 0$; this condition is the same as in the absence of viscosity}. This last result is in agreement with what was found in Chapter~XII, \S\S\,109 and 111\,(c).

While the system is stable for $\Sigma > 0$, the modes will be either damped-periodic or damped-aperiodic depending on the magnitude of $\Phi$, i.e.\ on the magnitude of the viscous dissipation.

\section{Extension of the variational principle to allow for compressibility}\label{sec:123}

The variational principle derived in the preceding section can be extended to allow for compressibility. Incompressibility simplifies the analysis by $\boldsymbol{\xi}$ being solenoidal; this latter property cannot be used if the fluid is compressible. On examining the details of the derivation in \S\,122, we find that the solenoidal property of $\boldsymbol{\xi}$ was used only in two places: first, in reducing the integral
\begin{equation}
\label{eq:14-43}
\int \frac{\partial \varpi^{(n)}}{\partial x_i} \xi_i^{(m)} \, dV
\end{equation}
to the simple surface integral in~(\ref{eq:14-18}); and again in substituting from~(\ref{eq:14-10}) for $\mathbf{h}$ in the volume integral on the right-hand side of equation~(\ref{eq:14-18}) and writing it as on the left-hand side of equation~(\ref{eq:14-19}). The terms which have been neglected on these two accounts and which we must now include are
\begin{equation}
\label{eq:14-44}
I = -\int \varpi^{(n)} \frac{\partial \xi_i^{(m)}}{\partial x_i} \, dV
\end{equation}
and
\begin{equation}
\label{eq:14-45}
\mathit{II} = \frac{1}{4\pi} \int \frac{\partial \xi_i^{(n)}}{\partial x_k} \left( H_j \frac{\partial H_k}{\partial x_j} \xi_i^{(m)} - H_i H_j \frac{\partial \xi_i^{(m)}}{\partial x_j} \right) dV.
\end{equation}

The reduction of the surface integral which remains in~(\ref{eq:14-18}), as well as the others, does not depend on the solenoidal character of $\boldsymbol{\xi}$ and can be carried out as before. The net result, then, is that on the right-hand side of equation~(\ref{eq:14-28}) we must now include the terms $I$ and $\mathit{II}$ with the others already present.

Considering $I$, we rewrite it in the form
\begin{equation}
\label{eq:14-46}
I = -\int \left(\delta p^{(n)} + \frac{H_j h_j^{(n)}}{4\pi}\right) \frac{\partial \xi_i^{(m)}}{\partial x_i} \, dV,
\end{equation}
where $\delta p^{(n)}$ is the perturbation in the material pressure. We must now express~(\ref{eq:14-46}) like the others, in terms of $\boldsymbol{\xi}$; to accomplish this, a physical hypothesis is needed to relate $\delta p$ with the displacement $\boldsymbol{\xi}$; and it is customary to suppose that the variations of density and pressure take place adiabatically as we follow fluid elements with their motions. In the Eulerian framework this means
\begin{equation}
\label{eq:14-47}
\frac{\partial}{\partial t} \delta p + u_j \frac{\partial p}{\partial x_j} = \gamma \frac{p}{\rho}\left(\frac{\partial}{\partial t} \delta \rho + u_j \frac{\partial \rho}{\partial x_j}\right),
\end{equation}
where $\gamma$ denotes the ratio of the specific heats and $p$ is the pressure on the unperturbed configuration. In addition to~(\ref{eq:14-47}) we also have the equation of continuity
\begin{equation}
\label{eq:14-48}
\frac{\partial}{\partial t} \delta \rho + u_j \frac{\partial \rho}{\partial x_j} = -\rho \frac{\partial u_j}{\partial x_j}.
\end{equation}
In terms of $\boldsymbol{\xi}$ equations~(\ref{eq:14-47}) and~(\ref{eq:14-48}) can be integrated to give
\begin{equation}
\label{eq:14-49}
\delta p = -\xi_k \frac{\partial p}{\partial x_k} - \gamma p \frac{\partial \xi_k}{\partial x_k}
\end{equation}
and
\begin{equation}
\label{eq:14-50}
\delta \rho = -\xi_k \frac{\partial \rho}{\partial x_k} - \rho \frac{\partial \xi_k}{\partial x_k},
\end{equation}
where $c$ denotes the velocity of sound. Inserting accordingly for $\delta p^{(n)}$ in~(\ref{eq:14-46}), we have
\begin{equation}
\label{eq:14-51}
\int \left(\xi_k^{(n)} \frac{\partial p}{\partial x_k} + c^2 \rho \frac{\partial \xi_k^{(n)}}{\partial x_k}\right) \frac{\partial \xi_i^{(m)}}{\partial x_i} \, dV + \int c^2 \rho \frac{\partial \xi_i^{(m)}}{\partial x_i} \frac{\partial \xi_k^{(n)}}{\partial x_k} \, dV + \int \xi_i^{(n)} \frac{\partial \xi_i^{(m)}}{\partial x_i} \frac{\partial p}{\partial x_k} \, dV.
\end{equation}
Similarly substituting for $\mathbf{h}^{(n)}$ from equation~(\ref{eq:14-10}) in the other term in~(\ref{eq:14-46}), we obtain
\begin{align}
\label{eq:14-52}
&-\frac{1}{4\pi} \int H_k \left(H_j \frac{\partial \xi_i^{(n)}}{\partial x_j} - H_k \frac{\partial H_j}{\partial x_j} \xi_i^{(n)}\right) \frac{\partial \xi_i^{(m)}}{\partial x_i} \, dV \notag \\
&\quad = -\frac{1}{4\pi} \int |\mathbf{H}|^2 \frac{\partial \xi_i^{(n)}}{\partial x_i} \frac{\partial \xi_i^{(m)}}{\partial x_i} \, dV + \int \xi_i^{(n)} \frac{\partial \xi_i^{(m)}}{\partial x_i} \frac{\partial}{8\pi} |\mathbf{H}|^2 \, dV \notag \\
&\qquad\qquad - \frac{1}{4\pi} \int H_j H_k \frac{\partial \xi_i^{(n)}}{\partial x_k} \frac{\partial \xi_i^{(m)}}{\partial x_i} \, dV.
\end{align}

We thus obtain
\begin{align}
\label{eq:14-53}
I &= \int \left(c^2 \rho + \frac{|\mathbf{H}|^2}{4\pi}\right) \frac{\partial \xi_i^{(n)}}{\partial x_i} \frac{\partial \xi_k^{(m)}}{\partial x_k} \, dV + \int \xi_i^{(n)} \frac{\partial \xi_k^{(m)}}{\partial x_k} \frac{\partial}{\partial x_i} \left(p + \frac{|\mathbf{H}|^2}{8\pi}\right) dV \notag \\
&\quad - \frac{1}{4\pi} \int H_j H_k \frac{\partial \xi_i^{(n)}}{\partial x_k} \frac{\partial \xi_i^{(m)}}{\partial x_i} \, dV.
\end{align}
Now making use of equation~(\ref{eq:14-1}) governing the stationary state, we can rewrite equation~(\ref{eq:14-53}) in the manner (cf.\ equation~(\ref{eq:14-45}))
\begin{equation}
\label{eq:14-54}
I = \int \left(c^2 \rho + \frac{|\mathbf{H}|^2}{4\pi}\right) \frac{\partial \xi_i^{(n)}}{\partial x_i} \frac{\partial \xi_k^{(m)}}{\partial x_k} \, dV + \frac{1}{4\pi} \int \left(H_k \frac{\partial H_j}{\partial x_j} \xi_i^{(n)} - H_j H_k \frac{\partial \xi_i^{(n)}}{\partial x_k}\right) \frac{\partial \xi_i^{(m)}}{\partial x_i} \, dV.
\end{equation}

Now adding the terms $I$ and $\mathit{II}$ given by equations~(\ref{eq:14-45}) and~(\ref{eq:14-54}) to the right-hand side of equation~(\ref{eq:14-28}), we finally obtain
\begin{align}
\label{eq:14-55}
(\sigma^{(n)})^2 \int_{V^{(\mathrm{in})}} \rho \xi_i^{(n)} \xi_i^{(m)} \, dV &= \int_S (\mathbf{N}^{(0)} \cdot \boldsymbol{\xi}^{(n)})(\mathbf{N}^{(0)} \cdot \boldsymbol{\xi}^{(m)})[\mathbf{N}^{(0)} \cdot \operatorname{grad}\, \Delta_S(\Pi)] \, dS \notag \\
&\quad + \frac{1}{4\pi} \int \left[\left(H_j \frac{\partial \xi_i^{(n)}}{\partial x_j}\right)\left(H_k \frac{\partial \xi_i^{(m)}}{\partial x_k}\right) + 4\pi \xi_i^{(n)} \xi_i^{(m)} \frac{\partial^2 \Pi}{\partial x_k \partial x_k}\right] dV \notag \\
&\quad + \int \left(c^2 \rho + \frac{|\mathbf{H}|^2}{4\pi}\right) \left|\frac{\partial \xi_i}{\partial x_i}\right|^2 dV + \frac{1}{4\pi} \int_{V^{(\mathrm{ex})}} h_i^{(n)} h_i^{(m)} \, dV \notag \\
&\quad + \frac{1}{4\pi} \int \left\{ H_k \frac{\partial H_j}{\partial x_k} \xi_i^{(n)} \frac{\partial \xi_i^{(m)}}{\partial x_j} - H_j H_k \frac{\partial \xi_i^{(n)}}{\partial x_k} \frac{\partial \xi_i^{(m)}}{\partial x_j} \right\} dV \notag \\
&\quad - \frac{1}{2\pi} \int H_j H_k \frac{\partial \xi_i^{(n)}}{\partial x_k} \left(\xi_i^{(m)} \frac{\partial \xi_i}{\partial x_j} + \frac{\partial \xi_i^{(m)}}{\partial x_j}\right) dV.
\end{align}

We observe that, as in the incompressible case, the expression on the right-hand side of equation~(\ref{eq:14-55}) is symmetrical in $\lambda$ and $\mu$; and, therefore,
\begin{equation}
\label{eq:14-56}
\int_{V^{(\mathrm{in})}} \rho \xi_i^{(n)} \xi_i^{(m)} \, dV = 0 \quad (\lambda \neq \mu).
\end{equation}

This clearly establishes the self-adjoint character of the problem. Also by writing the complex conjugates of $\xi_i^{(n)}$ and $h_i^{(n)}$ in place of $\xi_i^{(m)}$ and $h_i^{(m)}$ and suppressing the superscripts, we obtain
\begin{align}
\label{eq:14-57}
\sigma^2 \int_{V^{(\mathrm{in})}} \rho |\boldsymbol{\xi}|^2 \, dV &= \int_S |\mathbf{N}^{(0)} \cdot \boldsymbol{\xi}|^2 [\mathbf{N}^{(0)} \cdot \operatorname{grad}\, \Delta_S(\Pi)] \, dS \notag \\
&\quad + \frac{1}{4\pi} \int \left[\left|H_j \frac{\partial \xi_i}{\partial x_j}\right|^2 + 4\pi \xi_i \xi_i^* \frac{\partial^2 \Pi}{\partial x_k \partial x_k}\right] dV \notag \\
&\quad + \int \left(c^2 \rho + \frac{|\mathbf{H}|^2}{4\pi}\right) \left|\frac{\partial \xi_i}{\partial x_i}\right|^2 dV + \frac{1}{4\pi} \int_{V^{(\mathrm{ex})}} |\mathbf{h}|^2 \, dV \notag \\
&\quad + \frac{1}{2\pi} \int \left\{H_k \frac{\partial H_j}{\partial x_k} \xi_i^{(n)} \frac{\partial \xi_i^{(m)*}}{\partial x_j} - H_j H_k \frac{\partial \xi_i}{\partial x_k} \frac{\partial \xi_i^*}{\partial x_j}\right\} dV.
\end{align}

From this equation it follows that $\sigma^2$ \emph{is necessarily real and that a necessary and sufficient condition for stability is that the right-hand side of equation~\emph{(\ref{eq:14-57})} be positive}.

From the self-adjoint character of the problem it is apparent that equation~(\ref{eq:14-57}) provides the basis for a variational treatment of the problem.

\section*{BIBLIOGRAPHICAL NOTES}

The variational principles discussed in this chapter are derived from:
\begin{enumerate}
\item I.~B.~\textsc{Bernstein}, E.~A.~\textsc{Frieman}, M.~D.~\textsc{Kruskal}, and R.~M.~\textsc{Kulsrud}, `An energy principle for hydromagnetic stability problems', \emph{Proc.\ Roy.\ Soc.\ (London)} A, \textbf{244}, 17--40 (1958). \cite{Bernstein1958}
\item K.~\textsc{Hain}, R.~\textsc{L\"ust}, and A.~\textsc{Schl\"uter}, `Zur Stabilit\"at eines Plasmas', \emph{Z.\ f.\ Naturforschung}, \textbf{12a}, 833--41 (1957). \cite{Hain1957}
\item A.~\textsc{Hare}, `The effect of viscosity on the stability of incompressible magnetohydrodynamic systems', \emph{Phil.\ Mag.}, Ser.~8, \textbf{48}, 1305--10 (1959). \cite{Hare1959}
\item T.~G.~\textsc{Cowling}, `Magneto-hydrodynamic oscillations of a rotating fluid globe', \emph{Proc.\ Roy.\ Soc.\ (London)} A, \textbf{233}, 319--22 (1955). \cite{Cowling1955}
\end{enumerate}
